\documentclass[11pt]{article}

\usepackage[margin=1.1in]{geometry}
\usepackage{amsmath,amssymb,amsthm}
\usepackage{mathtools}

\newtheorem{lemma}{Lemma}
\newtheorem{theorem}{Theorem}
\newtheorem{remark}{Remark}

\newcommand{\R}{\mathbb{R}}
\newcommand{\cp}{\mathrm{cp}}
\newcommand{\cpbar}{\overline{\mathrm{cp}}}
\newcommand{\dist}{\operatorname{dist}}
\newcommand{\intr}{\operatorname{int}}
\newcommand{\norm}[1]{\lVert #1 \rVert}
\newcommand{\ip}[2]{\langle #1, #2 \rangle}

\title{Consistency and stability of the $\cp-\cpbar$ conormal estimator\\
  \large (\texttt{angle2d.m} / \texttt{angle3D.m})}
\author{}
\date{\today}

\begin{document}
\maketitle

\begin{abstract}
We prove that the rotation-estimation method implemented in
\texttt{angle2d.m} and \texttt{angle3D.m} --- which averages normalized
vectors $\cp(x)-\cpbar(x)$ over grid points whose closest point lies on
the boundary of a surface (or at a branch/singular point), and builds a
rotation matrix from the averaged direction --- is first-order consistent
with the true outward conormal and the exact continuous rotation, and is
stable: the constructed matrix lies exactly in $SO(n)$, the
direction-to-rotation map is $1$-Lipschitz, and all normalizations are
$O(1)$-conditioned with degenerate inputs rejected. For straight
segments and planar patches the direction estimate is exact.
\end{abstract}

\section{Setup}

Let $S\subset\R^n$ ($n=2,3$) be a $C^2$ embedded manifold-with-boundary
of positive reach $r_0>0$ in the region of interest, with boundary
$\Gamma=\partial S$ (also of reach $\ge r_0$). For $q\in\Gamma$ write
$T_qS$ for the tangent space of $S$, $N_qS$ for its orthogonal
complement in $\R^n$, and $T_q\Gamma\subset T_qS$ for the tangent of the
boundary (in 2D, $\Gamma$ is a set of endpoints and $T_q\Gamma=\{0\}$).
Let $\nu(q)\in T_qS$ denote the \emph{unit outward conormal}: the vector
in $T_qS$ orthogonal to $T_q\Gamma$ pointing out of $S$ through
$\Gamma$; in 2D this is the outgoing unit tangent of the curve at the
endpoint. Let $\kappa$ bound the relevant curvatures (second fundamental
form of $S$, geodesic curvature of $\Gamma$).

Let $\cp=P_S$ be the closest-point projection, well defined and $C^1$ on
the tube $\{\dist(\cdot,S)<r_0\}$. Following Macdonald, Brandman and
Ruuth (2011), define the double projection
\[
  \cpbar(x) \;=\; \cp\bigl(2\,\cp(x)-x\bigr).
\]
At each grid point $x$ the code computes
\[
  v(x) \;=\; \cp(x) - \cpbar(x),
\]
keeps the points whose closest point lies on $\Gamma$ (the
\texttt{bdy} mask, or proximity of $\cp(x)$ to a specified singular
point within \texttt{singularTol}), discards vectors with
$\norm{v}\le\texttt{vectorTol}$, normalizes each survivor, averages the
unit vectors, renormalizes the average to obtain the direction estimate
(\texttt{tangent} in 2D, \texttt{conormal} in 3D), and constructs a
rotation matrix $R$ from it (directly in 2D; via Rodrigues' formula
about a prescribed axis in 3D).

Two standard facts about $P_S$ inside the tube are used throughout:
\begin{itemize}
\item[(P1)] If $\cp(x)=p\in\intr S$, then $x-p\in N_pS$: the residual is
  purely $S$-normal.
\item[(P2)] The set of points projecting to $q\in\Gamma$ is the normal
  cone
  \[
    \mathcal N_q=\bigl\{\,q+b\,\nu(q)+\eta \;:\; b\ge 0,\ \eta\in
    N_qS\,\bigr\}\cap\{\text{tube}\},
  \]
  where $b=\ip{x-q}{\nu(q)}\ge 0$ is the \emph{outward overhang} of $x$
  past the boundary.
\end{itemize}

\section{The geometric core}

\begin{lemma}[Interior annihilation]\label{lem:interior}
If $\cp(x)=p\in\intr S$, then $v(x)=0$ exactly.
\end{lemma}

\begin{proof}
By (P1), $x-p\in N_pS$. The reflected point is
$\bar x = 2p-x = p-(x-p)$, so $\bar x - p = -(x-p)\in N_pS$: $\bar x$
lies on the same normal line through $p$, within the reach. Every point
of that segment has closest point $p$ by definition of the projection,
so $\cpbar(x)=\cp(\bar x)=p$ and $v(x)=p-p=0$.
\end{proof}

Lemma~\ref{lem:interior} justifies the \texttt{vectorTol} filter:
interior and degenerate samples contribute a provably zero vector and
are removed; only genuine boundary points inform the estimate.

\begin{lemma}[Boundary expansion]\label{lem:boundary}
Let $\cp(x)=q\in\Gamma$, so by \textup{(P2)} $x-q = b\,\nu(q)+\eta$ with
$b\ge 0$ and $\eta\in N_qS$. Then
\[
  v(x) \;=\; b\,\nu(q) \;+\; O\!\bigl(\kappa\,(b^2+b\norm{\eta})\bigr).
\]
In particular, if $S$ is flat near $q$ (a straight segment or planar
patch, $\kappa=0$), then $v(x)=b\,\nu(q)$ exactly.
\end{lemma}

\begin{proof}
We have $\bar x = 2q-x = q - b\,\nu - \eta$. Split $\bar x - q$ into its
$S$-tangential part $-b\,\nu$ (pointing \emph{into} $S$, since
$b\ge 0$) and $S$-normal part $-\eta$. Let $\sigma$ be the unit-speed
geodesic in $S$ with $\sigma(0)=q$, $\sigma'(0)=-\nu$; then
$\sigma(b) = q - b\,\nu + O(\kappa b^2)$, and
\[
  \bar x - \sigma(b) = -\eta + O(\kappa b^2),
\]
which is $S$-normal up to $O(\kappa b\norm{\eta})$ (the normal spaces
along $\sigma$ tilt at rate $O(\kappa)$). For a competing foot point
$\sigma(b)+\tau$ with $\tau$ tangential,
\[
  \norm{\bar x-\sigma(b)-\tau}^2
  = \norm{\bar x-\sigma(b)}^2 + \norm{\tau}^2
    - 2\ip{\bar x-\sigma(b)}{\tau},
\]
and the cross term is
$O\bigl(\kappa(b^2+b\norm{\eta})\,\norm{\tau}\bigr)$ because
$\bar x-\sigma(b)$ is nearly normal. Minimizing over $\tau$ forces
$\norm{\tau}=O(\kappa(b^2+b\norm{\eta}))$, hence
\[
  \cpbar(x) = \cp(\bar x)
  = q - b\,\nu + O\bigl(\kappa(b^2+b\norm{\eta})\bigr),
\]
and $v(x)=q-\cpbar(x)=b\,\nu+O(\kappa(b^2+b\norm{\eta}))$. When
$\kappa=0$ the geodesic is a straight line in $S$, all error terms
vanish, and $v(x)=b\,\nu$ exactly for every $b\ge0$.
\end{proof}

\begin{remark}
The exact case is the situation of the V-shaped example
(\texttt{example\_v\_shaped.m}): each branch is a straight segment, so
\texttt{angle2d} recovers the outgoing branch tangent with zero
discretization error, up to floating point.
\end{remark}

\section{Consistency}

\begin{theorem}[Consistency]\label{thm:consistency}
Let $x_1,\dots,x_m$ be the retained candidate points, with
$\cp(x_i)=q_i\in\Gamma$, $\norm{x_i-q_i}\le Ch$ (a band of width
$O(h)$), and overhang bounded below, $b_i\ge c\,h$ for some $c>0$
(automatic for a standard band; vacuous when $\kappa=0$). Let
$\delta=\max_i\norm{q_i-q_0}$ measure the clustering of the footpoints
around the target boundary point $q_0$ (enforced by
\texttt{singularTol}, or $\delta\to0$ as $h\to0$ under the \texttt{bdy}
mask). Then the returned direction satisfies
\[
  \mathrm{conormal}
  \;=\; \nu(q_0) + O(\kappa h + \kappa\delta)
  \;\xrightarrow[h,\delta\to0]{}\; \nu(q_0),
\]
and the constructed rotation satisfies
$\norm{R-R_{\mathrm{exact}}} = O(\kappa h + \kappa\delta)$. When
$\kappa=0$ the estimates are exact.
\end{theorem}

\begin{proof}
By Lemma~\ref{lem:boundary}, since normalization removes the magnitude
$b_i$,
\[
  u_i \coloneqq \frac{v(x_i)}{\norm{v(x_i)}}
  = \nu(q_i) + O\!\Bigl(\frac{\kappa(b_i^2+b_i h)}{b_i}\Bigr)
  = \nu(q_i) + O(\kappa h),
\]
using $b_i\ge ch$ and $\norm{\eta_i}\le Ch$. Smoothness of $\Gamma$
gives $\nu(q_i)=\nu(q_0)+O(\kappa\delta)$. Averaging the clustered unit
vectors,
\[
  \bar a = \frac1m\sum_{i=1}^m u_i
  = \nu(q_0) + O(\kappa h + \kappa\delta),
  \qquad
  \norm{\bar a} = 1 - O\bigl((\kappa h+\kappa\delta)^2\bigr),
\]
so $\norm{\bar a}$ is bounded away from zero and the renormalized
average satisfies the stated bound.

For the rotation: in 2D,
$R=\begin{psmallmatrix} t_1 & -t_2\\ t_2 & t_1\end{psmallmatrix}$ is a
smooth ($1$-Lipschitz, Theorem~\ref{thm:stability}(b)) function of the
unit direction $t$. In 3D, with unit axis $u$ and $s,t$ the normalized
projections of the estimated conormal and the target direction onto
$u^\perp$, the code sets $\cos\theta=\ip{s}{t}$,
$\sin\theta=\ip{u}{s\times t}$ and applies Rodrigues' formula; $R$ is a
smooth function of $(\theta,u)$ wherever the projections exceed the
tolerance. In both cases $\norm{R-R_{\mathrm{exact}}}$ is bounded by the
direction error. Finally, composed branch rotations of the form
$R_{A\to B}=R_B\,R_\pi R_A^{\top}$ inherit the same bound, since for
orthogonal factors $\norm{AB-A'B'}\le\norm{A-A'}+\norm{B-B'}$ in
operator norm.
\end{proof}

\section{Stability}

\begin{theorem}[Stability]\label{thm:stability}
The method has the following stability properties.
\begin{itemize}
\item[(a)] \textbf{Exact rotations.} For any input direction the
  constructed matrix lies exactly in $SO(n)$. In 2D,
  $R^{\top}R=(t_1^2+t_2^2)I=I$ and $\det R=t_1^2+t_2^2=1$ because $t$ is
  normalized. In 3D, with $s,t$ unit and orthogonal to $u$, one has
  $s\times t=(\sin\varphi)\,u$, hence
  $\cos^2\theta+\sin^2\theta=\cos^2\varphi+\sin^2\varphi=1$ exactly, and
  with $K$ the skew matrix of the \emph{unit} axis ($K^{\top}=-K$,
  $K^3=-K$), Rodrigues'
  $R=I+\sin\theta\,K+(1-\cos\theta)K^2$ satisfies $R^{\top}R=I$ and
  $\det R=1$ algebraically. The method never emits a non-rotation.
\item[(b)] \textbf{$1$-Lipschitz direction-to-rotation map.} In 2D,
  $t\mapsto R(t)$ is the standard isometry $S^1\cong SO(2)$: a direction
  error of angle $\delta$ gives
  $\norm{R(t')-R(t)}=2\lvert\sin(\delta/2)\rvert\le\delta$. In 3D,
  $\partial R/\partial\theta$ is a rotation generator of unit operator
  norm on $u^\perp$, so an angle error $\delta$ yields
  $\norm{\Delta R}\le\delta$; the projections onto $u^\perp$ are
  bounded linear maps. Direction error is not amplified.
\item[(c)] \textbf{$O(1)$-conditioned normalizations.} Each per-point
  normalization $u_i=v_i/\norm{v_i}$ is performed only when
  $\norm{v_i}>\texttt{vectorTol}$, so the relative amplification of
  floating-point error in $v_i$ is $O(1)$. Averaging \emph{unit}
  vectors reduces independent per-point noise like $1/\sqrt m$, and the
  final division uses $\norm{\bar a}\ge\tfrac12$ (clustering,
  Theorem~\ref{thm:consistency}).
\item[(d)] \textbf{Rejection of ill-posed inputs.} If no valid vectors
  remain, if the unit vectors cancel
  ($\norm{\bar a}\le\texttt{vectorTol}$: a symmetric configuration with
  no defined direction), or if the conormal or target is parallel to
  the axis (projected norms below tolerance: the in-plane angle is
  undefined), the code returns \texttt{NaN} rather than an arbitrary
  bounded rotation.
\item[(e)] \textbf{Scale invariance.} The tolerance
  $\texttt{tol}=100\,\varepsilon\cdot\max(1,\max\lvert\text{data}\rvert)$
  is homogeneous of degree one in the coordinates, so the retained set
  and all threshold tests are invariant under a change of units.
\end{itemize}
\end{theorem}

\begin{proof}
Each item is proved in its statement; (a) and (b) are algebraic
identities, (c) follows from the tolerance guards and
Theorem~\ref{thm:consistency}, and (d), (e) are read off the
implementation.
\end{proof}

\section{Conclusion and hypotheses}

By Lemmas~\ref{lem:interior}--\ref{lem:boundary}, $\cp-\cpbar$ vanishes
identically over the interior and equals $b\,\nu+O(\kappa h^2)$ over the
boundary band; normalization removes the magnitude, so the averaged,
renormalized direction converges to the true outward conormal at rate
$O(\kappa h)$ (exactly for straight or planar branches,
Theorem~\ref{thm:consistency}). The rotation constructors map this
direction into $SO(n)$ exactly and $1$-Lipschitzly
(Theorem~\ref{thm:stability}), so the estimated rotation converges at
$O(h)$ with unit sensitivity, and degenerate inputs return
\texttt{NaN}. The method is therefore consistent and stable. Two
hypotheses deserve emphasis:

\begin{enumerate}
\item \emph{Positive reach / resolved geometry.} All arguments require
  $h\ll r_0$ so that $\cp$ is single-valued on the band. Near a genuine
  cusp, or where two sheets approach within $O(h)$ (e.g.\ the
  piriform-cusp and two-sphere-intersection examples), the local reach
  shrinks and the constant in $O(\kappa h)$ degrades; the mechanism is
  unchanged, but the rate is meaningful only once $h$ resolves the
  local reach.
\item \emph{Overhang lower bound $b_i\ge ch$.} Needed only for curved
  boundaries, to control the per-point direction error
  $O(\kappa h^2/b_i)$; it holds for a standard computational band and
  is irrelevant when $\kappa=0$.
\end{enumerate}

\begin{thebibliography}{9}
\bibitem{MBR2011}
C.~B. Macdonald, J.~Brandman, S.~J. Ruuth,
\emph{Solving eigenvalue problems on curved surfaces using the closest
point method},
J. Comput. Phys. 230 (2011), 7944--7956.
\end{thebibliography}

\end{document}
